\SourceLeaf{351}%
gel, welches jenem als ganze Krümmung entspricht, zu der ganzen Oberfläche der Hülfskugel. Allgemein wird der Überschuss der Winkel eines Polygons von \(n\) Seiten, wenn diese kürzeste Linien sind, über \(2n-4\) Rechte, der ganzen Krümmung des Polygons gleich sein.

Die allgemeinen in der Abhandlung entwickelten Untersuchungen werden am Schluss derselben noch auf die Theorie der durch kürzeste Linien gebildeten Dreiecke angewandt, wovon wir hier nur ein paar Haupttheoreme anführen. Sind \(a,b,c\) die Seiten eines solchen Dreiecks (die als Grössen der ersten Ordnung betrachtet werden); \(A,B,C\) die gegenüberstehenden Winkel; \(\alpha,\beta,\gamma\) die Krümmungsmaasse in den Winkelpunkten; \(\sigma\) der Flächeninhalt des Dreiecks, so ist, bis auf Grössen der vierten Ordnung, \(\tfrac13(\alpha+\beta+\gamma)\sigma\) der Überschuss der Summe \(A+B+C\) über zwei Rechte. Ferner sind, mit derselben Genauigkeit, die Winkel eines ebenen geradlinigen Dreiecks, dessen Seiten \(a,b,c\) sind, der Ordnung nach
\[
\begin{aligned}
&A-\tfrac1{12}(2\alpha+\beta+\gamma)\sigma\\
&B-\tfrac1{12}(\alpha+2\beta+\gamma)\sigma\\
&C-\tfrac1{12}(\alpha+\beta+2\gamma)\sigma
\end{aligned}
\]
Man sieht sogleich, dass das letzte Theorem eine Generalisirung des bekannten von \textsc{Legendre} zuerst aufgestellten ist, nach welchem man, bis auf Grössen der vierten Ordnung, die Winkel des geradlinigen Dreiecks erhält, wenn man die Winkel des sphärischen jeden um den dritten Theil des sphärischen Excesses vermindert. Auf einer nichtsphärischen Fläche muss man also den Winkeln ungleiche Reductionen beifügen, und die Ungleichheit ist allgemein zu reden eine Grösse der dritten Ordnung; wenn jedoch die ganze Fläche nur wenig von der Kugelgestalt abweicht, so involvirt jene noch ausserdem einen Factor von der Ordnung der Abweichung von der Kugelgestalt. Es ist unstreitig für die höhere Geodaesie wichtig, dass man im Stande ist, die Ungleichheiten jener Reductionen zu berechnen, und dadurch die volle Ueberzeugung zu erhalten, dass sie für alle messbaren Dreiecke auf der Oberfläche der Erde als ganz unmerklich zu betrachten sind. So. finden sich z.\,B. in dem grössten Dreiecke der von dem Verf. ausgeführten Triangulirung, dessen grösste Seite fast 15 geographische Meilen lang ist, und in welchem der Ueberschuss der Summe der drei Winkel über zwei Rechte fast 15 Secunden beträgt, die drei Reductionen der Winkel auf die Win%
